\chapter{Архитектурно-техническое решение задачи}\label{ch:5}

При проведении структурного системного анализа была создана модель "'as is"', которая показала, что значительная часть работы выполняется вручную. Мастера автосервиса тратят большое количество времени на общение с клиентом в мессенджере, анализ фотографий и получение информации, формирование предварительной оценки. Такой подход требует лишних временных затрат и увеличивает нагрузку на сотрудников. 

Чтобы решить описанную выше проблему, рассмотрим внедрение интеллектуальной системы в текущий процесс на BPMN-диаграмме "to be". (Рисунок 11). В диаграмме значительная часть операций текущего процесса передаётся новой автоматизированной системе.
Функционал, предоставляемый системой для пользователя, должен включать:

\begin{enumerate}[label=\arabic*.]
    \item Начало взаимодействия с системой через мессенджер. Пользователь имеет возможность написать боту и выбрать режим работы системы
    \item Выбор способа указания повреждений автомобиля. Система должна обеспечивать возможность выбрать один из двух способов указания повреждений:
    \begin{itemize}
        \item Указать повреждения и их типы через интерфейс бота
        \item Отправить фотографию повреждённого автомобиля, которую автоматически проанализирует система
    \end{itemize}
    \item Если пользователь выбрал автоматический режим:
    \begin{itemize}
        \item Система должна автоматически выполнить обработку и анализ фотографии с использованием методов машинного обучения: определить видимые детали автомобиля, имеющиеся на них повреждения и сформировать аннотированное изображение с выделенными масками деталей и повреждений.
        Должны обнаруживаться следующие детали:
        \begin{itemize}
            \item Дверь
		    \item Переднее крыло
		    \item Заднее крыло
		    \item Бампер
		    \item Багажник
		    \item Капот
		    \item Фара
		    \item Переднее стекло
		    \item Заднее стекло
		    \item Боковое стекло
		    \item Шины и в целом колёса
        \end{itemize}
        Должны обнаруживаться следующие повреждения:
        \begin{itemize}
            \item Царапина
		    \item Вмятина
		    \item Отвалившийся кусок краски
		    \item Ржавчина
		    \item Трещина
		    \item Разрывы/разрезы/дырки
	        \item Разбитые стекло или фары
		    \item Спущенные шины и любые проблемы с колёсами
        \end{itemize}
        \item Пользователь должен иметь возможность посмотреть результат автоматического анализа перед переходом к этапу расчёта стоимости ремонта. 
        \item Пользователь должен иметь возможность отредактировать результат анализа через интерфейс бота, если он не согласен с результатом, выданным моделью:
        \begin{itemize}
            \item Изменить тип повреждения
            \item Изменить принадлежность повреждения к детали автомобиля
            \item Удалить ошибочно найденное повреждение
            \item Добавить новое необнаруженное повреждение
        \end{itemize}
    \end{itemize}
    \item После проверки пользователем всех параметров повреждений система должна предоставить ему возможность подтвердить итоговый список повреждений автомобиля, чтобы перейти к расчёту.
    \item После подтверждения от пользователя система должна приступить к автоматическому расчёту ориентировочных стоимости и длительности работ в автосервисе согласно заданным параметрам, сопоставив их с правилами расчёта автосервиса (Таблицы 5 и 6)
    \item Пользователь должен получить сообщение с результатами расчёта, включающими стоимость и длительность ремонтных работ. 
\end{enumerate}


\begin{figure}[H]
\centering
\begin{minipage}{0.90\textwidth}
    \centering
    \includegraphics[scale=0.4, angle=90]{BPMN_to_be}
\end{minipage}\hspace{-2em}%
\begin{minipage}{0.06\textwidth}
    \centering
    \rotatebox{90}{%
        \parbox{10cm}{\caption{BPMN-диаграмма TO BE}}%
    }
\end{minipage}
\end{figure}

% \begin{table}[htbp]
% \centering
% \caption{Стоимость ремонта (тыс. руб.)}

% \begin{adjustbox}{max width=\textwidth}
% \begin{tabular}{lcccc}
% \hline
% Деталь & Царапина (полировка) & Царапина (покраска) & Вмятина (рихтовка + покраска) & Замена детали \\
% \hline

% Переднее крыло & 1 & 10--18 & 23--30 & 20 \\
% Заднее крыло & 1 & 10--18 & 23--30 & 75--100 \\
% Дверь & 1 & 10--18 & 23--30 & 20 \\
% Крышка багажника & 1 & 10--18 & 23--30 & 20 \\

% Порог & 1 & 10--18 & 23--30 & 50--55 \\
% Крыша & 1 & 10--18 & 23--30 & 75--100 \\

% Капот & 1 & 10--18 & 30--35 & 28--30 \\

% Бампер & 1 & 10--18 & 3--5 & 18 \\

% Лобовое стекло & --- & --- & --- & 5+ \\
% Боковое стекло & --- & --- & --- & 3 \\
% Фара & --- & --- & --- & 3 \\

% \hline
% \end{tabular}
% \end{adjustbox}

% \end{table}


% \begin{table}[htbp]
% \centering
% \caption{Длительность ремонта}

% \begin{adjustbox}{max width=\textwidth}
% \begin{tabular}{lcccc}
% \hline
% Деталь & Царапина (полировка) & Царапина (покраска) & Вмятина (рихтовка + покраска) & Замена детали \\
% \hline

% Переднее крыло & 1 час & 1 день & 2--3 дня & 1 день \\
% Заднее крыло & 1 час & 1 день & 2--3 дня & 5 дней \\
% Дверь & 1 час & 1 день & 2--3 дня & 1.5--2 дня \\
% Крышка багажника & 1 час & 1 день & 2--3 дня & 1.5--2 дня \\

% Порог & 1 час & 1 день & 2--3 дня & 1 день \\
% Крыша & 1 час & 1 день & 2--3 дня & 5 дней \\

% Капот & 1 час & 1 день & 2 дня & 2 дня \\

% Бампер & 1 час & 1 день & 1--2 дня & 1.5 дня \\

% Лобовое стекло & --- & --- & --- & 3 часа (оставить на день) \\
% Боковое стекло & --- & --- & --- & 1 день \\
% Фара & --- & --- & --- & 0.5 дня \\

% \hline
% \end{tabular}
% \end{adjustbox}

% \end{table}



\begin{table}[htbp]
\centering
\caption{Стоимость ремонта, тыс. руб.}
\begin{adjustbox}{max width=\textwidth}
\begin{tabular}{|l|c|c|c|c|}
\hline
Деталь & \begin{tabular}[c]{@{}c@{}}Царапина\\ (полировка)\end{tabular} & 
\begin{tabular}[c]{@{}c@{}}Царапина\\ (покраска)\end{tabular} & 
\begin{tabular}[c]{@{}c@{}}Вмятина\\ (рихтовка + покраска)\end{tabular} & 
Замена детали \\
\hline

Переднее крыло & 1 & 10–18 & 23–30 & 20 \\
Заднее крыло & 1 & 10–18 & 23–30 & 75–100 \\
Дверь & 1 & 10–18 & 23–30 & 20 \\
Крышка багажника & 1 & 10–18 & 23–30 & 20 \\

% Порог & 1 & 10–18 & 23–30 & 50–55 \\
Крыша & 1 & 10–18 & 23–30 & 75–100 \\

Капот & 1 & 10–18 & 30–35 & 28–30 \\

Бампер & 1 & 10–18 & 3–5 & 18 \\

Лобовое стекло & — & — & — & 5-10 \\
Боковое стекло & — & — & — & 3 \\
Фара & — & — & — & 3 \\

\hline
\end{tabular}
\end{adjustbox}
\end{table}

\begin{table}[htbp]
\centering
\caption{Длительность ремонта}
\begin{adjustbox}{max width=\textwidth}
\begin{tabular}{|l|c|c|c|c|}
\hline
Деталь & \begin{tabular}[c]{@{}c@{}}Царапина\\ (полировка)\end{tabular} & 
\begin{tabular}[c]{@{}c@{}}Царапина\\ (покраска)\end{tabular} & 
\begin{tabular}[c]{@{}c@{}}Вмятина\\ (рихтовка + покраска)\end{tabular} & 
Замена детали \\
\hline

Переднее крыло & 1 час & 1 день & 2–3 дня & 1 день \\
Заднее крыло & 1 час & 1 день & 2–3 дня & 5 дней \\
Дверь & 1 час & 1 день & 2–3 дня & 1.5–2 дня \\
Крышка багажника & 1 час & 1 день & 2–3 дня & 1.5–2 дня \\

% Порог & 1 час & 1 день & 2–3 дня & 1 день \\
Крыша & 1 час & 1 день & 2–3 дня & 5 дней \\

Капот & 1 час & 1 день & 2 дня & 2 дня \\

Бампер & 1 час & 1 день & 1–2 дня & 1.5 дня \\

Лобовое стекло & — & — & — & 1 день \\
Боковое стекло & — & — & — & 1 день \\
Фара & — & — & — & 0.5 дня \\

\hline
\end{tabular}
\end{adjustbox}
\end{table}




После выполнения всех шагов выше клиент сможет проще принять решение о пользовании услугами автосервиса, и сделать лишь финальный звонок сотруднику автосервиса (о времени визита / согласовании цены).

Предлагаемое решение, встроенное в текущий процесс, имеет несколько преимуществ:

\begin{enumerate}[label=\arabic*.]
    \item При поступлении обращения через канал мессенджера можно полностью передать управление системе - клиенту достаточно будет отправить фотографию или самостоятельно указать все имеющиеся повреждения.
    \item Сотруднику автосервиса не требуется проводить долгое общение в мессенджере для формирования ответа по стоимости - в системе клиент это может сделать самостоятельно, и, благодаря удобному интерфейсу, быстро задать параметры и получить ответ.
    \item Помимо полной автоматизации процесса при поступлении обращений через мессенджеры, можно предлагать клиентам пользоваться системой и при поступлении телефонных звонков - это также будет кратно экономить время.
\end{enumerate}

Для реализации функционала необходимо разработать приложение, которое будет обеспечивать описанные сценарии. Архитектура и математическая модель системы будут рассмотрены в следующем разделе.

